\documentclass[11pt]{article}

\usepackage[a4paper,margin=0.85in]{geometry}
\usepackage{setspace}
\usepackage{graphicx}
\usepackage{array}
\usepackage{parskip}
\usepackage{ragged2e}

\pagestyle{empty}
\setstretch{1.05}

\begin{document}

\begin{flushright}
Zhongli Wang \\
+49 160 9541 6015 \\
nkwangzhongli@gmail.com \\
Astastrasse 22, 33617 Bielefeld \\
Germany \\
\end{flushright}

\vspace{0.1cm}

\noindent
Search Committee \\
Department of Engineering and Applied Sciences \\
Columbia University \\
Infosys Topaz Columbia University Enterprise AI Center \\
New York, NY \\

\vspace{0.1cm}

\noindent
\textbf{Application for Associate Research Scientist in Enterprise AI} \\

\vspace{0.05cm}

\noindent Dear Members of the Search Committee,

I am currently completing my Ph.D.\ in Economics at Bielefeld University and Universit\`a Cattolica del Sacro Cuore, with planned completion in September 2026. My research combines data analysis, machine learning, statistical modeling, agent-based simulation, financial and economic data systems, and emerging work on agentic research workflows. I am very interested in the Associate Research Scientist position at the Infosys Topaz Columbia University Enterprise AI Center because the Center's focus on responsible, human-centered, and enterprise-oriented AI closely matches my developing research agenda: building reliable AI-enabled systems that support complex research, decision-making, and organizational workflows.

My strongest fit with this position lies in the intersection of knowledge-grounded models, data management, enterprise automation, and responsible AI adoption. During my doctoral studies, I have developed end-to-end analytical workflows for complex economic and financial questions, including data collection, cleaning, validation, feature engineering, model construction, robustness checks, and the preparation of structured research outputs. This experience has shaped my interest in AI systems that do not merely generate outputs, but support transparent, auditable, and reusable workflows. In enterprise settings, where data are heterogeneous, decisions are high-stakes, and users need traceable evidence, I believe the design of reliable knowledge-grounded AI systems is both technically important and practically urgent.

A central direction of my recent work is the development of agentic workflows for economic research. In my project \emph{Agentic Workflows for Economic Research}, I explore how large-language-model agents can be integrated into empirical research pipelines, including document collection, classification, structured extraction, workflow documentation, and research assistance. My interest is not simply in using LLMs as general-purpose chat tools, but in designing domain-specific systems with clear task boundaries, human review, source traceability, and reproducible outputs. This perspective is directly relevant to the Center's emphasis on enterprise automation, knowledge-grounded models, AI reliability, and the translation of research prototypes into practical platforms.

My doctoral dissertation, \emph{Strategic Behavior and Economic Dynamics in ICT Markets: An Agent-Based Framework}, studies how firms and heterogeneous users interact under network effects, congestion constraints, imperfect information, and strategic investment decisions. Methodologically, this work required me to formalize behavioral decision rules, build scalable simulation pipelines, evaluate feedback mechanisms, and conduct counterfactual analysis. These skills are relevant to enterprise AI research because the deployment of AI systems in organizations also involves strategic behavior, feedback loops, information frictions, and unintended consequences. I am particularly interested in studying how AI agents and human users interact in complex workflows, how reliability can be evaluated under realistic behavioral responses, and how governance mechanisms can reduce risks while preserving productivity gains.

A second line of my work connects machine learning with financial and organizational risk. In collaborative research on \emph{Polytope Fraud Theory}, I use supervised machine-learning methods to detect structured anomalies in financial statements, with a focus on firms targeted by activist short sellers. This work involved constructing firm-level features, handling heterogeneous financial indicators, evaluating model performance through recall, F1 score, predictive accuracy, and feature importance, and interpreting algorithmic signals in terms of financial misconduct, investor protection, and market transparency. This experience has strengthened my ability to move from a substantive problem to data construction, model development, validation, and interpretation. It has also made me attentive to the reliability, auditability, and governance of AI-supported decision systems.

I can contribute to the Center's research activities in three ways. First, I can help design and prototype domain-specific agentic applications for research and enterprise workflows, especially systems that combine document processing, structured data extraction, classification, retrieval, and human-in-the-loop review. Second, I can contribute to empirical evaluation frameworks for AI reliability, including robustness testing, error analysis, workflow comparison, and the assessment of whether AI-supported tools improve the quality, speed, and transparency of knowledge work. Third, I can support collaborative publications, open-source demonstrations, and applied research outputs that connect technical rigor with practical business impact.

My technical background includes machine learning, statistical modeling, simulation, data processing, and reproducible research workflows. I have strong programming experience in Python, including pandas, NumPy, scikit-learn, XGBoost, and related tools for data analysis, modeling, and experimentation. I also have experience with Stata, R, SQL, Linux-based research environments, and structured project workflows. I am continuing to deepen my knowledge of LLM application frameworks, retrieval-augmented generation, agentic systems, evaluation methods, and responsible AI governance. I would be highly motivated to further develop these skills in collaboration with Columbia faculty, Ph.D.\ students, and Infosys engineers.

Beyond technical research, I bring strong experience in academic writing, presentation, and collaborative research communication. My Ph.D.\ training has required me to write research papers, prepare conference presentations, document methodological decisions, and communicate complex quantitative findings to different audiences. I would be glad to mentor students and interns, contribute to exploratory research directions, assist with proposal preparation, and participate actively in conferences, academic collaborations, and industry-facing forums. I am also open to relocating to the New York City area for this position.

I am particularly attracted to the Infosys Topaz Columbia University Enterprise AI Center because it combines academic research with large-scale enterprise AI challenges. The opportunity to work at the interface of original research, prototype development, responsible AI governance, and industry deployment is especially appealing to me. I would be excited to contribute to the Center's intellectual and practical mission by developing reliable, knowledge-grounded, and human-centered AI systems for complex enterprise and research environments.

Thank you so much for considering my application! I would welcome the opportunity to discuss how my background and abilities can contribute to the research.

\vspace{0.15cm}

\noindent Sincerely, \\[0.8cm]
\noindent Zhongli Wang

\end{document}
